% Página
\section{Criptografia do Lado do Cliente}

Neste contexto, \textbf{criptografia do lado do cliente} (também chamada de arquitetura de conhecimento zero) significa que os dados são cifrados e decifrados no dispositivo do usuário, e o serviço é projetado para não receber o texto plano, a Senha Mestra ou as chaves de dados em claro. O termo não se refere a uma prova criptográfica de conhecimento zero, que é uma técnica distinta para provar uma afirmação sem revelar o segredo usado na prova.

\subsection{Funcionamento da Criptografia do Lado do Cliente}

\begin{enumerate}
	\item Derivação da chave
		\begin{itemize}
			\item O cliente combina a Senha Mestra com um salt aleatório e usa uma KDF apropriada para derivar uma chave. O salt e os parâmetros da KDF podem ser armazenados junto aos dados cifrados; eles não precisam ser secretos.
		\end{itemize}
	\item Cifragem local
		\begin{itemize}
			\item O cliente cifra as credenciais e outros dados localmente antes de enviá-los ao serviço.
			\item Deve-se usar uma cifra autenticada (AEAD), como AES-GCM, com um nonce novo e único para cada operação sob a mesma chave. O texto cifrado, o nonce e a tag de autenticação são enviados e armazenados juntos.
		\end{itemize}
	\item Armazenamento no serviço
		\begin{itemize}
			\item O serviço armazena os dados cifrados e os metadados necessários para derivar a chave e decifrá-los, mas não deve receber as chaves em claro.
			\item O serviço ainda pode conhecer metadados não cifrados, como tamanhos, horários e informações da conta, e pode apagar, reter ou reverter dados armazenados.
		\end{itemize}
	\item Recuperação e decifragem
		\begin{itemize}
			\item O cliente baixa os dados cifrados, deriva novamente a chave com a Senha Mestra e o salt armazenado e verifica a tag AEAD antes de liberar o texto plano.
			\item Se a verificação falhar, a chave ou os parâmetros podem estar incorretos, ou os dados podem ter sido alterados; a falha, por si só, não identifica a causa.
		\end{itemize}
\end{enumerate}

Esse modelo reduz o acesso do provedor ao conteúdo, mas não torna os dados automaticamente seguros diante de qualquer comprometimento. Uma senha mestra fraca pode permitir tentativas de adivinhação offline contra os dados cifrados; um dispositivo comprometido pode capturar a senha ou o texto plano. Sem uma chave de recuperação independente, perder a Senha Mestra pode tornar os dados irrecuperáveis.

\pagebreak
